\documentclass[letterpaper]{article}

\usepackage[margin=1.5cm]{geometry}
\usepackage{amsmath,amsfonts}
\usepackage{url}
\usepackage{bm}

\newcommand{\Real}{\mathbb{R}}
\newcommand{\tr}{^\top}
\newcommand{\opt}{^\star}

\setlength{\parskip}{3mm plus1mm minus1mm}
\setlength{\parindent}{0mm}

\title{Problem Set 13}
\author{Foundations of Machine Learning}
\date{}

\begin{document}
\maketitle

\begin{itemize}
\item \textbf{Submission}: Turn in as a \texttt{PDF} and a \textbf{separate source code file} on \url{http://mycourses.unh.edu}. Scanned hand-written notes are encouraged.
\item \textbf{Questions}: See the class website for a link to office hours
\end{itemize}

% Covers: Overfitting and cross-validation (MML 8.2). Cross-validation exercises are drawn from ISL ch.~5.

\textbf{Relevant reading}: MML 8.2 (empirical risk minimization, regularization, cross-validation); ISL 5 (resampling and cross-validation). Problems marked \textbf{[code]} should be submitted as a program in a language of your choice (Julia or Python recommended). The cited exercise numbers are from ISL, 2nd edition (James et al., 2021).

\begin{enumerate}
\item ISL, Chapter 5, Conceptual exercises: problems 1, 2, 3.
\item Explain, in a short paragraph, the bias--variance trade-off and how it relates to overfitting. Give one concrete example of an under-fit model and one of an over-fit model for the same data.
\item \textbf{[code]} Fit polynomials of increasing degree to a small data set. Use $k$-fold cross-validation to select the degree, and plot training error versus validation error as a function of degree.
\item \textbf{Time spent.} Report approximately how many hours you spent completing this problem set (including the reading).
\end{enumerate}

\end{document}
